\subsection{Interpretation and limitations}
\label{subsec:rewrite-rsi-limitations}

The model supports firmer conditional conclusions about long-run growth
and income distribution than about takeoff dates. In the finite-upper-bound
regimes characterized in Proposition~\ref{prop:rewrite-equilibrium-regimes},
once AI efficiency approaches $\overline B$, the limiting growth rates,
interest rate and income shares do not depend on $\chi>0$.
They depend on substitutability, the efficiency bound and the other
production and macroeconomic parameters. Research productivity instead
affects how quickly the economy approaches those limits and how resources
are allocated along the way. Its uncertain value makes the reported
transition dates particularly uncertain. This distinction does not imply
that $\sigma$ or $\overline B$ is well identified: both describe technological
scenarios, not estimated probabilities.

The central scenario and the slow-transition sensitivity are informed by
data, not joint empirical fits of the
current AI economy. At unit elasticity, AI-industry revenue is
$(1-\alpha)\omega_X=6.7\%$ of output, independently of $K_0$, $B_0$ and
$\chi$. That share is not fitted to the size estimates in
\citet{korinekmckelvey2026}. The expenditure proxy measures historical US
training compute, not worldwide autonomous RSI. Matching its level also
does not reproduce its recent growth. Nor do the exercises match
expenditure and price changes simultaneously. The half-price-decline
target in the central scenario is illustrative, not an estimate of the contribution
of RSI to observed price changes.

A more aggressive calibration, retained in the replication files, assigns
the full rounded $80\%$ observed price decline to RSI and gives
$\chi=36.830702$. At $\sigma=1.50$, it implies initial output-per-person
growth of $23.97\%$, research spending of $9.32\%$ of output, developer
net profit of $-6.02\%$ and gross investment of $-6.09\%$ of output.
Although it passes the numerical equilibrium checks, it requires large
developer losses and immediate liquidation of capital. I do not use it
as a main quantitative scenario because these initial reallocations are
implausibly large. The example exposes the model's adjustment assumptions:
RSI appears unexpectedly, compute expands without installation delays,
capital is reversible, and owners can finance losses against future profit.
Lower $\chi$ moderates the response but does not add the capacity costs or
complementary investments in \citet{erdiletal2025gate} and
\citet{brynjolfssonrocksyverson2021}. Appendix~\ref{app:rewrite-algorithm}
links to the preserved exercise and its replication instructions.

Permanent monopoly and the absence of human research further restrict
the interpretation of prices and developer profit. The income shares
describe payments to labor, capital and the AI developer. Without
heterogeneous asset ownership or the creation of new human tasks, they
do not establish how gains are distributed across households or how
political power changes.
